\documentclass[aspectratio=169,10pt]{beamer}
\newcommand{\reporoot}{../..}
% Shared preamble of the main deck and the topic decks.
% Define \reporoot (the path from the .tex file to the repository root)
% before \input-ing this file.
\usetheme[progressbar=frametitle,numbering=fraction,sectionpage=progressbar]{metropolis}

\usepackage{amsmath,amssymb}
\usepackage{booktabs}
\usepackage{graphicx}
\usepackage{tikz}
\usetikzlibrary{arrows.meta,positioning,decorations.pathreplacing}
\usepackage[authoryear,round]{natbib}
\usepackage{xcolor}

\graphicspath{{\reporoot/figures/}}

% Every number on the slides is produced by experiments/*.py and read from here.
\input{\reporoot/results/e1_markov_book.tex}
\input{\reporoot/results/e2_best_queues.tex}
\input{\reporoot/results/e3_queue_reactive.tex}
\input{\reporoot/results/e4_hawkes.tex}
\input{\reporoot/results/e5_signature.tex}
\input{\reporoot/results/e6_propagator.tex}
\input{\reporoot/results/e7_market_making.tex}
\input{\reporoot/results/e8_ofi.tex}
\input{\reporoot/results/e9_queue_value.tex}
\input{\reporoot/results/e10_markouts.tex}
\input{\reporoot/results/e11_glft_exact.tex}
\input{\reporoot/results/e12_signal.tex}
\input{\reporoot/results/e13_lvr.tex}
\input{\reporoot/results/e14_mlofi.tex}
\input{\reporoot/results/e15_rate_estimation.tex}
\input{\reporoot/results/e16_queue_reactive_price.tex}
\input{\reporoot/results/e17_hawkes_pitfalls.tex}
\input{\reporoot/results/e18_propagator_estimation.tex}
\input{\reporoot/results/e19_latent_book.tex}
\input{\reporoot/results/e20_microprice.tex}
\input{\reporoot/results/e21_fill_estimation.tex}
\input{\reporoot/results/e22_glosten_milgrom.tex}
\input{\reporoot/results/e23_cdl_on_book.tex}

\newcommand{\E}{\mathbb{E}}
\newcommand{\Var}{\operatorname{Var}}
\newcommand{\OFI}{\mathrm{OFI}}
\newcommand{\dd}{\mathrm{d}}
\newcommand{\src}[1]{\par\vfill{\scriptsize\color{gray}#1}}
\definecolor{accent}{HTML}{0C5DA5}
\definecolor{warn}{HTML}{C0392B}
\newcommand{\hl}[1]{\textcolor{accent}{\textbf{#1}}}

\setbeamerfont{footnote}{size=\tiny}
\setbeamerfont{caption}{size=\scriptsize}

\author{Yaohui Shu}
\institute{Stony Brook University}
\date{October 2026}

% Topic decks: a references frame that lists only the works cited in the deck.
\newcommand{\deckreferences}{%
  \begin{frame}[allowframebreaks]{References}
    \tiny
    \renewcommand{\bibsection}{}
    \bibliographystyle{plainnat}
    \bibliography{\reporoot/slides/refs}
  \end{frame}}

% Topic decks: the standard "where this fits" footer on the title frame.
\newcommand{\deckseries}{Deep dive in the series \emph{Stochastic Models of the Limit Order Book}}

% Topic decks: \decknumber{n} puts the series line under the date on the title page.
\newcommand{\decknumber}[1]{\institute{Stony Brook University\hfill{\scriptsize\color{gray}Deep dive #1 of 9}}}


\title{Price Impact}
\subtitle{Propagators, the square-root law, and the latent order book}
\decknumber{5}

\begin{document}

\maketitle

\begin{frame}{The question}
  \small
  \begin{columns}[T]
    \begin{column}{0.55\textwidth}
      Trades move prices. \hl{How much, for how long, and why is it concave in size?}
      \begin{itemize}
        \item Order signs have long memory, yet prices are close to unpredictable. Impact must decay just fast enough to offset the memory.
        \item A metaorder's impact grows like the square root of its size, across markets and decades.
        \item For an LP, impact is adverse selection seen from the other side: the price move that follows a fill.
      \end{itemize}
    \end{column}
    \begin{column}{0.42\textwidth}
      \begin{alertblock}{In this deck}
        \small
        \begin{enumerate}
          \item The propagator model and its critical decay.
          \item Estimating the propagator from trades.
          \item The latent order book: an exact crossover from linear to square-root impact (derived here).
          \item Evidence and open questions.
        \end{enumerate}
      \end{alertblock}
    \end{column}
  \end{columns}
\end{frame}

\begin{frame}{Two facts that must coexist}
  \small
  \begin{columns}[T]
    \begin{column}{0.5\textwidth}
      \textbf{Long memory of order signs.} With trade signs $\varepsilon_n = \pm1$ in trade time $n$, the autocorrelation decays as a power law,
      \[ C_\varepsilon(\ell) = \E[\varepsilon_n\varepsilon_{n+\ell}] \sim \ell^{-\gamma_\varepsilon},\qquad 0 < \gamma_\varepsilon < 1. \]
      On the London Stock Exchange $\gamma_\varepsilon \approx 0.6$ \citep{lillo2004long}. Large orders are split, so signs persist for thousands of trades.
    \end{column}
    \begin{column}{0.46\textwidth}
      \textbf{Square-root impact.} A metaorder of $Q$ shares, executed over hours, moves the price on average by
      \[ I(Q) \approx Y\sigma_D\sqrt{Q/V_D}, \]
      with daily volatility $\sigma_D$ (fraction of price), daily volume $V_D$ (shares) and $Y$ of order one \citep{toth2011anomalous}. More generally $I \propto Q^{\delta_Q}$ with $\delta_Q$ between 0.4 and 0.7.
    \end{column}
  \end{columns}
  \vspace{0.4em}
  If each trade had a permanent impact, memory would make prices trend. If impact were linear in size, splitting would not matter. Both are false in data.
\end{frame}

\begin{frame}{The propagator model}
  \small
  The price just before trade $n$, in ticks \citep{bouchaud2004fluctuations}:
  \[ P_n = \sum_{m<n} G(n - m)\,\varepsilon_m + \text{noise}, \qquad G(\ell) = \Gamma_0\,\ell^{-\zeta}, \]
  with $\Gamma_0$ the impact one trade later (ticks) and $\zeta$ the decay exponent. For $C_\varepsilon \sim \ell^{-\gamma_\varepsilon}$, the variance per trade of price changes over $\ell$ trades scales as $\mathcal{D}(\ell) = \Var(P_{n+\ell} - P_n)/\ell \sim \ell^{\,1 - \gamma_\varepsilon - 2\zeta}$, so prices diffuse only at $\zeta_c = (1 - \gamma_\varepsilon)/2$. Faster decay gives mean reversion, slower decay gives trends.
  \vspace{0.2em}
  \centering
  \includegraphics[width=0.56\textwidth]{fig6_propagator.png}

  \raggedright\footnotesize Signs from fractional Gaussian noise, $\gamma_\varepsilon = \EsixGamma$, $2^{20}$ trades. Fitted slopes $\EsixSlopeLow$, $\EsixSlopeCrit$, $\EsixSlopeHigh$; theory $\EsixTheoryLow$, $\EsixTheoryCrit$, $\EsixTheoryHigh$.
\end{frame}

\begin{frame}{Estimating the propagator}
  \small
  \begin{columns}[T]
    \begin{column}{0.52\textwidth}
      Differencing the model, the price change at trade $n$ is
      \[ \Delta P_n = \sum_{k\geq 0} g_k\,\varepsilon_{n-k} + \xi_n, \]
      with $g_0 = G(1)$, $g_k = G(k+1) - G(k)$, and noise $\xi_n$ (ticks) uncorrelated with past signs. Multiply by $\varepsilon_{n-j}$ and average:
      \[ \sum_{k=0}^{K-1} C_\varepsilon(|j - k|)\,g_k = S(j) = \E[\Delta P_n\,\varepsilon_{n-j}], \]
      a Toeplitz system for $g_0, \ldots, g_{K-1}$; then $G(\ell) = \sum_{k<\ell} g_k$ \citep{bouchaud2004fluctuations,eisler2012price}.
    \end{column}
    \begin{column}{0.45\textwidth}
      \textbf{Two warnings.}
      \begin{itemize}
        \item \hl{Truncation bias.} Lags beyond $K$ are left out, and long-memory signs make them correlated with the lags kept. $\hat G$ is biased near $\ell = K$. Choose $K$ much larger than the lags you care about.
        \item \hl{Response is not impact.} The observable $\mathcal{R}(\ell) = \E[\varepsilon_n(P_{n+\ell} - P_n)]$ mixes the propagator with the memory of future signs. It can rise while $G$ falls.
      \end{itemize}
      \textbf{Uncertainty.} Block bootstrap over contiguous blocks, longer than the sign memory.
    \end{column}
  \end{columns}
\end{frame}

\begin{frame}{Estimation in practice}
  \centering
  \includegraphics[width=0.86\textwidth]{fig18_propagator_estimation.png}
  \begin{columns}[T]
    \begin{column}{0.5\textwidth}
      \footnotesize
      $\EeighteenN$ trades, $\gamma_\varepsilon = 0.5$, $G(\ell) = \ell^{-1/4}$ (critical), plus price noise of 1 tick per trade. (a) With $K = 100$, $\hat G(100) = \EeighteenGhundredShort$ against the true $\EeighteenGtrue$, and $\hat\zeta = \EeighteenZetaShort$ (fit on lags 2 to 100). With $K = 2000$: $\hat G(100) = \EeighteenGhundredLong$, $\hat\zeta = \EeighteenZetaLong$.
    \end{column}
    \begin{column}{0.46\textwidth}
      \footnotesize
      (b) The response function stays near one tick ($\mathcal{R}(1) = \EeighteenRone$, maximum $\EeighteenRmax$) while the propagator falls to $0.15$: reading $\mathcal{R}$ as impact decay would conclude that impact is permanent. Errors: block bootstrap, 32 blocks.
    \end{column}
  \end{columns}
\end{frame}

\begin{frame}{The latent order book}
  \small
  The visible book holds a small fraction of daily volume, so it cannot explain impact over hours. \citet{donier2015fully} model the \hl{latent} book: intentions to trade that are revealed as the price approaches.
  \vspace{0.2em}
  \begin{columns}[T]
    \begin{column}{0.52\textwidth}
      Latent bid and ask densities $b(x,t)$, $a(x,t)$ (shares per tick) at price $x$ (ticks) diffuse with coefficient $D$ (ticks$^2$ s$^{-1}$) and annihilate on contact. With $\phi = b - a$ and a buy metaorder at rate $m(t)$ (shares s$^{-1}$):
      \[ \partial_t\phi = D\,\partial_x^2\phi + m(t)\,\delta_D(x - y_t), \qquad \phi(y_t, t) = 0, \]
      where $y_t$ is the price and $\delta_D$ the Dirac delta. Far from the price $\phi(x, 0) = -Lx$, with slope $L$ (shares tick$^{-2}$).
    \end{column}
    \begin{column}{0.45\textwidth}
      \textbf{Units first.} The only dimensionless group made of $m$, $L$, $D$ is
      \[ u = \frac{m}{LD}, \]
      execution rate over the rate at which diffusion brings liquidity to the price. For constant $m$ there is no time scale, so the price path is self-similar:
      \[ y_t = \sqrt{Dt}\;F(u). \]
      Simplest version: no deposition, no cancellation (infinite memory).
    \end{column}
  \end{columns}
\end{frame}

\begin{frame}{An exact crossover from linear to square-root impact}
  \small
  The equation is linear in $\phi$, so $\phi = -Lx + m\int_0^t K_D(x - y_s, t - s)\,\dd s$ with the heat kernel $K_D(x,\tau) = e^{-x^2/4D\tau}/\sqrt{4\pi D\tau}$. Set $\phi(y_t, t) = 0$ with $y_s = A\sqrt{s}$, substitute $s = tw$, $w = 1 - v^2$, and write $a = A/\sqrt{D}$:
  \[ \boxed{\;a = \frac{u}{\sqrt\pi}\int_0^1 \exp\!\Big(-\frac{a^2}{4}\,\frac{1 - \sqrt{1 - v^2}}{1 + \sqrt{1 - v^2}}\Big)\dd v\;} \qquad (F(u) = a). \]
  The left side increases in $a$ and the right side decreases, so the root is unique.
  \vspace{0.2em}
  \begin{columns}[T]
    \begin{column}{0.5\textwidth}
      \textbf{Slow execution} ($u \ll 1$): $a \to u/\sqrt\pi$, so for $Q = mT$ shares over $T$ s
      \[ I = \frac{Q}{L\sqrt{\pi D T}} \quad \text{(linear in } Q). \]
    \end{column}
    \begin{column}{0.46\textwidth}
      \textbf{Fast execution} ($u \gg 1$): the integrand is concentrated near $v = 0$, $a \to \sqrt{2u}$, and
      \[ I = \sqrt{2Q/L} \quad \text{(square root, independent of } D, T). \]
    \end{column}
  \end{columns}
  \vspace{0.2em}
  Checks: units ($I$ in ticks in both limits); the fast limit equals the price at which a linear book $Lx$ holds $Q$ shares. Derived for this repository; \citet{donier2015fully} give the two regimes.
\end{frame}

\begin{frame}{Closed form against the PDE}
  \centering
  \includegraphics[width=0.88\textwidth]{fig19_latent_book.png}
  \footnotesize
  \begin{columns}[T]
    \begin{column}{0.48\textwidth}
      (a) Crank-Nicolson solution of the PDE on 16,001 nodes: within $\EnineteenMaxErr\%$ of the closed form up to $u = \EnineteenUmax$; halving the resolution doubles the error ($\EnineteenErrCoarse\%$), first-order convergence. (b) The local exponent falls from $\EnineteenExpoLow$ to $\EnineteenExpoHigh$; the asymptotes cross at $u = 2\pi$.
    \end{column}
    \begin{column}{0.48\textwidth}
      (c) After the metaorder, impact decays. In the linear regime the exact path is
      \[ y_t = \tfrac{m}{L\sqrt{\pi D}}\big(\sqrt{t} - \sqrt{t - T}\big), \]
      a propagator with exponent $1/2$ in calendar time. At $t = 5T$ the slow metaorder keeps $\EnineteenRemainSlow$ of its peak (theory $\EnineteenRemainLin$); the fast one keeps $\EnineteenRemainFast$.
    \end{column}
  \end{columns}
\end{frame}

\begin{frame}{What the data show}
  \small
  \begin{itemize}
    \item \textbf{Concave impact everywhere.} Reported exponents $\delta_Q$ range from 0.4 to 0.7 across studies. On about 700,000 US stock orders executed by Citigroup, \citet{almgren2005direct} find a 3/5 power law. On nearly 500,000 trades in futures, \citet{toth2011anomalous} find $\delta_Q \approx 0.5$ for small-tick and $\approx 0.6$ for large-tick contracts, for $Q/V_D$ from a few $10^{-4}$ to a few percent.
    \item \textbf{Sign memory.} Order signs on the LSE have autocorrelation decaying roughly as $\ell^{-0.6}$, offset by anticorrelated liquidity and volume \citep{lillo2004long}.
    \item \textbf{All events have impact.} Limit orders and cancellations move prices too; for large-tick stocks a model with permanent impact per event fits, for small-tick stocks history dependence matters \citep{eisler2012price}.
    \item \textbf{Consistency.} The latent book model gives square-root impact, contains the propagator model, and allows no price manipulation \citep{donier2015fully}. Exponential decay of impact is compatible only with linear impact \citep{gatheral2010no}.
  \end{itemize}
\end{frame}

\begin{frame}{Discussion and takeaways}
  \small
  \begin{columns}[T]
    \begin{column}{0.5\textwidth}
      \textbf{For a liquidity provider}
      \begin{itemize}
        \item The fill of a resting order is the start of someone's impact. A fill in the middle of a metaorder is followed by more of the same sign: long-horizon markouts measure this.
        \item Impact decays, so part of the adverse move reverts. The split between transient and permanent impact sets whether to hedge or wait.
        \item Do not estimate impact decay from the response function; use the propagator system with a long truncation lag.
      \end{itemize}
    \end{column}
    \begin{column}{0.46\textwidth}
      \textbf{Open questions}
      \begin{itemize}
        \item One mechanism for queue dynamics, sign memory and square-root impact.
        \item Inferring the latent book online from fills and visible depth.
        \item The latent book with finite memory: how much of impact becomes permanent, and on what time scale.
      \end{itemize}
      \textbf{Not tested here.} The square-root law is empirical; this repository solves one model of it, with no market data.
    \end{column}
  \end{columns}
\end{frame}

\begin{frame}[fragile]{Reproduce}
  \small
  Code: \texttt{impact.py} (propagator, estimation), \texttt{latent\_book.py}. Experiments: \texttt{e6\_propagator.py}, \texttt{e18\_propagator\_estimation.py}, \texttt{e19\_latent\_book.py}.
\begin{verbatim}
from stochastic_lob import latent_book as lb, impact as im
lb.impact_peak(Q=100.0, T=1.0, L=1.0, D=1.0)        # closed form, ticks
t, y = lb.simulate_metaorder(m=100.0, T=1.0)        # PDE check
G = im.estimate_propagator(signs, dprice, max_lag=2000)
\end{verbatim}
\end{frame}

\deckreferences
\end{document}
